\documentclass[11pt,a4paper]{article}

\usepackage[T1]{fontenc}
\usepackage{mlmodern}
\usepackage[margin=2.3cm]{geometry}
\usepackage{amsmath}
\usepackage{booktabs}
\usepackage{enumitem}
\usepackage{xcolor}
\usepackage{listings}
\usepackage[colorlinks=true,linkcolor=black,urlcolor=black,citecolor=black,filecolor=black]{hyperref}

\definecolor{codebg}{RGB}{247,247,249}
\definecolor{codekw}{RGB}{0,70,140}
\definecolor{codecm}{RGB}{110,110,110}
\definecolor{saybg}{RGB}{240,246,255}

\lstset{
  language=Python,
  basicstyle=\ttfamily\footnotesize,
  keywordstyle=\color{codekw}\bfseries,
  commentstyle=\color{codecm}\itshape,
  backgroundcolor=\color{codebg},
  frame=single, rulecolor=\color{black!15},
  breaklines=true, columns=fullflexible, keepspaces=true,
  showstringspaces=false, xleftmargin=4pt, xrightmargin=4pt
}
\setlist{itemsep=2pt,topsep=3pt}
\setlength{\parindent}{0pt}
\setlength{\parskip}{6pt}

\newcommand{\code}[1]{\texttt{#1}}
\newcommand{\showfile}[1]{\par\noindent\textbf{Show:} \code{#1}\par}

\begin{document}

\begin{center}
{\large Reinforcement Learning Lab \;$\cdot$\; BCSE432E \;$\cdot$\; L55 + L56 \;$\cdot$\; Dr.\ Rajesh Kumar}\\[8pt]
{\LARGE\bfseries Presentation Script}\\[4pt]
{\large Adaptive Curriculum and Self-Play for Generalizable Adversarial RL Agents}\\[6pt]
Aayush Shukla (23BPS1081) \quad Aditya Kumar (23BPS1163)
\end{center}

\noindent\rule{\textwidth}{0.4pt}

\textit{How to use this: read the normal text out loud. ``Show'' tells you which file to open
on screen. Code boxes are lines from our project; the comments on the right were added to
explain them.}

%=====================================================================
\section{Opening}
Our project is a game between two AI agents: a \textbf{Hunter} and an \textbf{Evader}.
The Evader has to find a key, open a locked door and escape through the exit. The Hunter has to
catch it before that happens.

Every game is played on a new, randomly generated dungeon. So the real question is not
``can the agent win on a map it has practised on'', but ``can it win on a map it has
\emph{never seen}''. We also ask whether smart training methods, called \textbf{curriculum
learning} and \textbf{self-play}, give better agents than simple training.

We wrote the learning algorithms, \textbf{PPO} and \textbf{DQN}, ourselves from scratch. We did not
use any ready-made reinforcement learning library.

%=====================================================================
\section{The Game}
The map is a $21\times21$ grid with rooms, corridors, one locked door and one key. Each turn both
agents move at the same time: up, down, left, right or stay. A game ends after at most 200 steps.

Each agent sees only a small $9\times9$ area around itself, and walls block its view, like a
torch in a dark dungeon. That makes the game hard, because the agents often do not know where
the other one is.

\textbf{Rewards}, which are the agent's ``score'': $+1$ for winning, $-1$ for losing, a tiny penalty
every step so it does not waste time, and a small bonus for getting closer to its goal.

\showfile{env/generator.py}
Makes the maps. Each map comes from a seed number, so map 42 is always the same map. It also
checks that every map can actually be solved.

\showfile{env/dungeon\_env.py}
The game rules: moves, catching, line of sight, and the rewards.

%=====================================================================
\section{What Reinforcement Learning Is}
Reinforcement learning is learning by trial and error. The agent tries actions, gets rewards,
and slowly learns which actions lead to more reward.

Our agent learns two things:
\begin{itemize}
  \item \textbf{The policy}: given what I see, what is the chance of each move? This is the agent's
        decision maker.
  \item \textbf{The value} $V(s)$: from this situation, how much reward do I expect to get? This
        is the agent's sense of ``am I winning right now?''
\end{itemize}
Both are small neural networks. The input is what the agent sees, 496 numbers. The policy
network outputs 5 probabilities, one per move. The value network outputs one number.

\showfile{rl/nets.py}
\begin{lstlisting}
class ActorCritic(nn.Module):
    def __init__(self, obs_dim, n_actions, hidden=(256, 256)):
        super().__init__()
        self.pi = mlp([obs_dim, *hidden, n_actions], out_gain=0.01)   # policy: 496 -> 256 -> 256 -> 5
        self.v = mlp([obs_dim, *hidden, 1], out_gain=1.0)            # value:  496 -> 256 -> 256 -> 1
\end{lstlisting}
``Actor'' is the policy that acts. ``Critic'' is the value network that judges how good the
situation is.

%=====================================================================
\section{PPO: Our Main Algorithm}

\subsection*{The idea in simple words}
\begin{enumerate}
  \item Let the agent play many games (32 games at once, 256 steps each).
  \item For every move, ask: \emph{was this move better or worse than usual?} That number is
        called the \textbf{advantage}.
  \item Make good moves more likely and bad moves less likely.
  \item \textbf{But do not change too much at once.} That last rule is what PPO adds.
\end{enumerate}

Why the last rule? If you change the policy too much in one go, the agent can suddenly forget
everything and start playing badly. PPO stops this by limiting each change to about 20\%.

\subsection*{The one equation to show}
\[
\text{PPO objective} \;=\; \text{average of}\;\; \min\Big(\,\text{ratio}\times A,\;\;
\text{clip}(\text{ratio},\,0.8,\,1.2)\times A\,\Big)
\]
How to say it:
\begin{itemize}
  \item \textbf{ratio} = how much more (or less) likely the new policy makes this move than the old
        one. 1 means no change.
  \item $A$ = the advantage: positive means a good move, negative means a bad move.
  \item \textbf{clip(ratio, 0.8, 1.2)} = we do not let the ratio go beyond $\pm20\%$.
  \item Taking the \textbf{min} means we always use the more careful of the two. So the agent can
        improve, but only in small, safe steps.
\end{itemize}

\showfile{rl/ppo.py  (PPO.update)}
\begin{lstlisting}
ratio = log_ratio.exp()                                     # new prob / old prob
pg_loss = -torch.min(ratio * a,
                     ratio.clamp(1 - c.clip, 1 + c.clip) * a).mean()   # the equation above
v_loss = 0.5 * (self.net.value(obs[idx]) - ret_t[idx]).pow(2).mean()   # train the critic
entropy = -(logp_all.exp() * logp_all).sum(-1).mean()                   # keep exploring
loss = pg_loss + c.vf_coef * v_loss - c.ent_coef * entropy

self.opt.zero_grad()
loss.backward()                                             # compute gradients
self.opt.step()                                             # gradient descent step (Adam)
\end{lstlisting}
Line by line:
\begin{itemize}
  \item \code{pg\_loss} is the PPO equation. The minus sign is because the optimiser
        \emph{minimises}, and we want to \emph{maximise} the objective.
  \item \code{v\_loss} teaches the critic to predict the real returns better.
  \item \code{entropy} rewards the agent for not being too sure too early, so it keeps
        trying new moves.
  \item \code{backward()} and \code{step()} are the gradient descent: they nudge the network
        weights in the direction that lowers the loss.
\end{itemize}

\subsection*{How we compute ``better or worse than usual'' (the advantage)}
We use a method called \textbf{GAE}. In simple words: for each step, compare what actually
happened (the reward plus the value of where we ended up) with what the critic expected. If
reality was better than expected, the move was good. GAE adds these surprises up over the next
few steps, with older ones counting less.

\showfile{rl/ppo.py  (RolloutBuffer.compute\_gae)}
\begin{lstlisting}
for t in reversed(range(self.T)):
    nonterm = 1.0 - self.dones[t]
    delta = self.rewards[t] + gamma * next_v * nonterm - self.values[t]   # surprise at step t
    gae = delta + gamma * lam * nonterm * gae                             # add up future surprises
    adv[t] = gae
    next_v = self.values[t]
\end{lstlisting}

%=====================================================================
\section{DQN: Our Second Algorithm}
DQN learns a different thing: a \textbf{Q-value} $Q(s,a)$ for every move, meaning ``how good
is it to make move $a$ in situation $s$''. Then it simply picks the move with the highest Q-value.

\subsection*{The second equation to show}
\[
\text{target} \;=\; \text{reward} \;+\; 0.99 \times \text{best Q-value of the next situation}
\]
How to say it: ``The value of a move is the reward I get now, plus the value of the best thing I
can do next.'' The $0.99$ means future rewards count slightly less than immediate ones. The network
is trained to make its Q-value match this target.

Two tricks make DQN stable:
\begin{itemize}
  \item \textbf{Replay buffer}: store past experiences and learn from random old samples, not
        just the latest ones.
  \item \textbf{Target network}: a frozen copy of the network that computes the target, updated
        only every 2000 steps, so the target does not keep moving.
\end{itemize}

\showfile{rl/dqn.py  (DQN.train\_step)}
\begin{lstlisting}
q = self.q(o).gather(-1, a.unsqueeze(-1)).squeeze(-1)     # current Q-value of the move we made
with torch.no_grad():
    a2 = self.q(o2).argmax(-1, keepdim=True)              # best next move (online network)
    q2 = self.target(o2).gather(-1, a2).squeeze(-1)       # its value (target network)
    y = r + c.gamma * (1 - d) * q2                        # the equation above
loss = F.smooth_l1_loss(q, y)                             # make Q close to the target
\end{lstlisting}
\textbf{Why PPO is our main algorithm}: in our game the opponent keeps changing while it learns, so
old stored experiences go out of date, which hurts DQN. Also, PPO's policy is naturally
random-ish, which turned out to matter a lot (Section~\ref{sec:finding}).

%=====================================================================
\section{How We Train Two Agents Against Each Other}
We compare three ways of training:
\begin{itemize}
  \item \textbf{Single}: one agent learns against a fixed, hand-written opponent. This is the
        simple baseline.
  \item \textbf{Independent}: both agents learn at the same time. The problem is that each one is
        chasing a moving target. The Hunter learns a trick, the Evader learns to counter it, the
        Hunter forgets its old skill, and they go round in circles.
  \item \textbf{Self-play}: each agent plays against a \textbf{pool of frozen old copies} of the
        other one. Because the old copies do not change, the agent has a stable target, and
        because it keeps meeting old opponents, it does not forget how to beat them.
\end{itemize}

\showfile{rl/pool.py}
\begin{lstlisting}
def sample(self):
    if self.rng.random() < self.p_latest:          # 50%: play the newest copy
        return self.items[-1]
    return self.items[int(self.rng.integers(len(self.items)))]   # 50%: any older copy
\end{lstlisting}
Every 50{,}000 steps we save a new frozen copy into the pool, which holds at most 20.

\showfile{training/train.py}
The main training loop. It chooses the mode with \code{--mode single / independent / selfplay}.

%=====================================================================
\section{Curriculum Learning}
At the start the agent knows nothing. On the full map it almost never wins, so it has almost
nothing to learn from. Curriculum learning is like school: start with easy lessons, move up when
you are ready.

\begin{center}
\begin{tabular}{ccccc}
\toprule
Level & Map size & Doors & Obstacles & Opponent plays randomly\\
\midrule
0 & $13\times13$ & 0 & none & 75\% of the time\\
1 & $15\times15$ & 1 & few  & 50\%\\
2 & $17\times17$ & 1 & few  & 25\%\\
3 & $21\times21$ & 1 & more & never (full skill)\\
\bottomrule
\end{tabular}
\end{center}
The agent moves to the next level when it wins at least 60\% of its recent games. When both agents
are learning, we instead check that at least 80\% of games end with a catch or an escape, and not
a timeout.

\showfile{rl/curriculum.py}
\begin{lstlisting}
advance = (at_level >= self.min_steps and full and self.rate() >= self.threshold) \
    or at_level >= self.max_steps            # good enough, or stayed too long
if advance:
    self.level += 1
\end{lstlisting}

%=====================================================================
\section{How We Test the Agents}
\begin{itemize}
  \item \textbf{Unseen maps}: we train on maps 0--999 and test on maps 20000--20499, which the
        agent has never seen. We also test on bigger $25\times25$ maps with two doors, which are
        harder than anything it trained on.
  \item \textbf{Generalization gap}: win rate on training maps minus win rate on test maps. If it is
        close to zero, the agent learned real skills and did not just memorise maps.
  \item \textbf{Exploitability}: we freeze a trained Hunter and train a brand-new Evader against it.
        If the new Evader learns to beat it quickly, the Hunter has a weakness. A strong Hunter
        takes much longer to beat.
  \item \textbf{Cross-play}: we play old and new versions against each other to check whether
        training keeps improving or goes round in circles.
\end{itemize}

\showfile{evaluation/generalization.py, evaluation/exploitability.py, evaluation/crossplay.py}

%=====================================================================
\section{Seeing Inside the Agent (Heatmaps)}
Win rate tells us \emph{if} the agent is good, not \emph{why}. So we put the agent on every square
of the map and ask its value network how good that square is. We colour the map with those values
and draw an arrow for the move it would pick.

You can read the agent's thinking from this: for example, the Hunter's value is high near the exit,
because that is where it can block the Evader.

\showfile{visualization/decision\_map.py, visualization/live.py (press H in the viewer)}

%=====================================================================
\section{Our Most Interesting Finding}
\label{sec:finding}
We tested the trained Hunter in two ways: always picking its single most likely move, or
choosing randomly according to its probabilities.

\begin{center}
\begin{tabular}{lcc}
\toprule
How the Hunter picks moves & Win rate on unseen maps & Games that time out\\
\midrule
Always the most likely move & 45\% & 30--39\%\\
Random, by its probabilities & 70\% & about 10\%\\
\bottomrule
\end{tabular}
\end{center}
Why? The Hunter cannot see far. When it loses sight of the Evader, many places look exactly the
same to it, so a fixed rule makes it walk left-right-left-right forever. A bit of randomness breaks
that loop. This matches known theory: when an agent cannot see everything, a random-ish policy can
be the best one.

%=====================================================================
\section{Closing}
To sum up: we built a dungeon chase game with random maps, wrote PPO and DQN ourselves, trained the
agents in three different ways with and without a curriculum, tested them on maps they had never
seen, measured how easy they are to beat, and made heatmaps to explain their decisions.

%=====================================================================
\section*{Where Is What (for quick reference)}
\begin{center}
\begin{tabular}{ll}
\toprule
\textbf{Topic} & \textbf{File}\\
\midrule
Map generator & \code{env/generator.py}\\
Game rules, rewards & \code{env/dungeon\_env.py}\\
Neural networks & \code{rl/nets.py}\\
PPO (update, GAE) & \code{rl/ppo.py}\\
DQN & \code{rl/dqn.py}\\
Self-play pool & \code{rl/pool.py}\\
Curriculum & \code{rl/curriculum.py}\\
Running 32 games in parallel & \code{rl/vec\_env.py}\\
Training (all modes) & \code{training/train.py}, \code{training/train\_dqn.py}\\
Hand-written opponents & \code{agents/scripted.py}\\
Testing on unseen maps & \code{evaluation/generalization.py}\\
Exploitability, cross-play & \code{evaluation/exploitability.py}, \code{evaluation/crossplay.py}\\
Heatmaps, live viewer & \code{visualization/decision\_map.py}, \code{visualization/live.py}\\
Results page & \code{visualization/dashboard.py} $\rightarrow$ \code{experiments/dashboard.html}\\
\bottomrule
\end{tabular}
\end{center}

\section*{If They Ask\dots{} (short answers)}
\begin{description}[style=nextline]
  \item[What is PPO?] A method that improves the agent's policy in small, safe steps, so it never
        changes too much at once.
  \item[What is DQN?] A method that learns how good each move is, the Q-value, and picks the best one.
  \item[Difference between V and Q?] $V$ = how good the situation is. $Q$ = how good a specific move
        is in that situation.
  \item[What is the critic for?] It tells the policy what ``normal'' looks like, so we can tell if a
        move was better or worse than usual.
  \item[Why self-play?] Training against a mix of old frozen opponents is stable and stops the agent
        forgetting.
  \item[Why curriculum?] On the full map early rewards are too rare. Easy maps first give the agent
        something to learn from.
  \item[Why not use the GPU?] Most of the time is spent running the game itself, which runs on the
        CPU. The network is tiny, so the GPU would save very little time.
  \item[What did you write yourselves?] PPO, DQN, the networks, the parallel game runner, self-play,
        the curriculum, evaluation and visualization. We used PyTorch only for the basic maths
        (tensors and gradients).
\end{description}

\end{document}
